% =====================================================================
% IBFE Tutorials: mathematical formulation
% =====================================================================
\documentclass[aspectratio=169,11pt]{beamer}

% Match the visual language of the SIAM Waves talk.
\usetheme{Boadilla}
\usecolortheme{seahorse}
\useinnertheme{rectangles}
\useoutertheme{infolines}

\setbeamercolor{structure}{fg=blue!55!black}
\setbeamercolor{title}{fg=white,bg=blue!55!black}
\setbeamercolor{frametitle}{fg=blue!55!black}
\setbeamercolor{block title}{bg=blue!12,fg=blue!55!black}
\setbeamercolor{block body}{bg=blue!4}
\setbeamertemplate{navigation symbols}{}
\setbeamerfont{frametitle}{series=\bfseries,size=\large}
\setbeamerfont{title}{series=\bfseries,size=\Large}

\usepackage[T1]{fontenc}
\usepackage[utf8]{inputenc}
\DeclareFontShape{T1}{cmss}{b}{n}{<->ssub*cmss/bx/n}{}
\usepackage{graphicx}
% These prefixes allow compilation from either the repository root or this folder.
\graphicspath{{./}{../}}
\usepackage{amsmath,amssymb,bm,mathtools}
\usepackage{booktabs}
\usepackage{xcolor}
\usepackage{tcolorbox}
\tcbuselibrary{skins}
\usepackage{hyperref}

\hypersetup{colorlinks=true,urlcolor=blue!55!black}

\newcommand{\hi}[1]{\textcolor{blue!55!black}{\textbf{#1}}}
\newcommand{\warn}[1]{\textcolor{red!70!black}{\textbf{#1}}}
\newcommand{\soft}[1]{\textcolor{gray!50!black}{#1}}

\newcommand{\xvec}{\bm{x}}
\newcommand{\Xvec}{\bm{X}}
\newcommand{\uvec}{\bm{u}}
\newcommand{\fvec}{\bm{f}}
\newcommand{\Fvec}{\bm{\mathcal{F}}}
\newcommand{\Gvec}{\bm{G}}
\newcommand{\Tvec}{\bm{T}}
\newcommand{\Vvec}{\bm{V}}
\newcommand{\Nvec}{\bm{N}}
\newcommand{\nvec}{\bm{n}}
\newcommand{\fzero}{\bm{f}_0}
\newcommand{\avec}{\bm{a}}
\newcommand{\mvec}{\bm{m}}
\newcommand{\Ften}{\bm{F}}
\newcommand{\Fcoef}{\bm{c}_{\!\mathcal F}}
\newcommand{\Pten}{\bm{P}}
\newcommand{\sigten}{\bm{\sigma}}
\newcommand{\Cten}{\bm{C}}
\newcommand{\Eten}{\bm{E}}
\newcommand{\Identity}{\bm{I}}
\newcommand{\Grad}{\nabla_{\!X}}
\newcommand{\Div}{\nabla_{\!X}\!\cdot}
\newcommand{\gradx}{\nabla_{\!x}}
\newcommand{\divx}{\nabla_{\!x}\!\cdot}

\newtcolorbox{punchbox}{
  enhanced,
  colback=blue!4,
  colframe=blue!55!black,
  boxrule=0.6pt,
  arc=2pt,
  left=4pt,right=4pt,top=3pt,bottom=3pt
}

\title[IBFE formulation]{IBFE mathematical formulation}
\subtitle{Strong form, weak form, and discrete coupling}
\author[Miller]{Laura A. Miller}
\institute[UArizona]{Department of Mathematics, University of Arizona}
\date{}

\begin{document}

% ----- 1. Title -----
\begin{frame}[plain]
  \begin{beamercolorbox}[sep=10pt,center]{title}
    {\bfseries\Large IBFE mathematical formulation}\\[3pt]
    {\normalsize Strong form, weak form, and discrete coupling}
  \end{beamercolorbox}
  \vspace{1.4em}
  \begin{columns}[T,onlytextwidth]
  \begin{column}{0.56\textwidth}
    \textbf{\large Laura A. Miller}\\[4pt]
    Department of Mathematics\\
    University of Arizona\\[14pt]
    {\small IBFE Tutorials}
  \end{column}
  \begin{column}{0.44\textwidth}
    \begin{punchbox}
    \centering
    \[
      W(\Ften)
      \ \longrightarrow\
      \Pten
      \ \longrightarrow\
      \Fvec_h
      \ \longrightarrow\
      \fvec_h
    \]
    \small Solid mechanics in reference coordinates, fluid dynamics on a Cartesian grid.
    \end{punchbox}
  \end{column}
  \end{columns}
  \vfill
  \centering{\tiny\color{black!55}AI tools were used to assist with formatting, organization, and clarification of the instructor's original notes.}
\end{frame}

\section{Overview and notation}

% ----- 2. Map -----
\begin{frame}{Roadmap from stress to coupled motion}
\begin{enumerate}\itemsep8pt
  \item Describe deformation by the motion \(\chi(\Xvec,t)\) and its gradient \(\Ften\).
  \item Convert a constitutive stress \(\Pten\) into a Lagrangian force by virtual work.
  \item Spread that force to the Eulerian fluid grid with a regularized delta kernel.
  \item Restrict fluid velocity back to the finite-element structure with an adjoint operator.
\end{enumerate}
\begin{punchbox}
\centering
\small The finite-element mesh represents geometry and mechanics. Quadrature points mediate the fluid--structure interaction.
\end{punchbox}
\end{frame}

% ----- 3. Notation bridge -----
\begin{frame}{Notation bridge from the solid-mechanics module}
\begin{columns}[T,onlytextwidth]
\begin{column}{0.52\textwidth}
\textbf{Kinematics and stress}
\[
  \Ften=\Grad\chi,
  \qquad
  J=\det\Ften,
  \qquad
  \Pten=\frac{\partial W}{\partial\Ften}.
\]
\begin{itemize}\itemsep5pt
  \item \(\Ften\): deformation gradient;
  \item \(\Pten\): first Piola--Kirchhoff stress;
  \item \(W\): energy per reference volume.
\end{itemize}
\end{column}
\begin{column}{0.48\textwidth}
\textbf{Force and velocity coupling}
\begin{itemize}\itemsep7pt
  \item \(\Fvec(\Xvec,t)\): Lagrangian force per reference volume;
  \item \(\fvec(\xvec,t)\): Eulerian force per current volume;
  \item \(\chi_t=\uvec\circ\chi\): no-slip motion.
\end{itemize}
\end{column}
\end{columns}
\begin{punchbox}
\small Here we write the reference and current solid regions as \(B_0\) and \(B_t\), reserving \(\Omega\) for the Eulerian computational domain.
\end{punchbox}
\end{frame}

\section{Kinematics and stress measures}

% ----- 4. Geometry -----
\begin{frame}{Reference and current configurations}
\begin{columns}[T,onlytextwidth]
\begin{column}{0.49\textwidth}
\[
  \chi(\cdot,t):B_0\longrightarrow B_t\subset\Omega,
  \qquad
  \xvec=\chi(\Xvec,t).
\]
\[
  \Ften(\Xvec,t)=\Grad\chi(\Xvec,t),
\]
\[
  J(\Xvec,t)=\det\Ften(\Xvec,t)>0.
\]
\begin{itemize}\itemsep5pt
  \item \(\Xvec\in B_0\): material coordinate.
  \item \(\xvec\in\Omega\): spatial coordinate.
  \item \(B_t=\chi(B_0,t)\): region occupied by the immersed solid.
\end{itemize}
\end{column}
\begin{column}{0.51\textwidth}
\centering
\includegraphics[width=\linewidth]{pictures/4-3/IBFE-sketch.png}
\end{column}
\end{columns}
\begin{punchbox}
\small IBFE keeps the structural description Lagrangian while momentum, viscosity, and incompressibility remain Eulerian.
\end{punchbox}
\end{frame}

% ----- 5. Stress -----
\begin{frame}{Fluid and structural stress measures}
\begin{columns}[T,onlytextwidth]
\begin{column}{0.55\textwidth}
\textbf{Newtonian fluid stress}
\[
  \sigten^f=-p\Identity
  +\mu\!\left(\gradx\uvec+(\gradx\uvec)^{\mathsf T}\right).
\]
\textbf{Structural stress transformation}
\[
  \sigten^s(\chi(\Xvec,t),t)
  =\frac{1}{J}\Pten(\Xvec,t)\Ften^{\mathsf T},
\]
\[
  \Pten\equiv\Pten^s=\Pten^p+\Pten^a,
  \qquad
  \Pten^p=\frac{\partial W}{\partial\Ften}.
\]
\end{column}
\begin{column}{0.45\textwidth}
\textbf{Inside and outside the solid}
\[
\sigten(\xvec,t)=
\begin{cases}
\sigten^f+\sigten^s,&\xvec\in B_t,\\[2pt]
\sigten^f,&\xvec\in\Omega\setminus B_t.
\end{cases}
\]
\small The superscript \(s\) denotes the total structural stress. In a passive calculation, \(\Pten^a=\bm{0}\).
\end{column}
\end{columns}
\begin{punchbox}
\small The Piola transform \(\Pten^s=J\sigten^s\Ften^{-\mathsf T}\) converts between current-area and reference-area stress measures.
\end{punchbox}
\end{frame}

% ----- 6. Assumptions -----
\begin{frame}{Classical neutrally buoyant IBFE assumptions}
\begin{columns}[T,onlytextwidth]
\begin{column}{0.49\textwidth}
\textbf{Fluid and solid properties}
\begin{itemize}\itemsep7pt
  \item viscous incompressible fluid;
  \item neutrally buoyant immersed solid;
  \item one Eulerian density \(\rho\);
  \item one fluid-like viscous stress \(\sigten^f\).
\end{itemize}
\end{column}
\begin{column}{0.51\textwidth}
\textbf{Coupling assumptions}
\begin{itemize}\itemsep7pt
  \item no slip and no penetration;
  \item the structure moves with \(\uvec\);
  \item solid inertia is carried by the Eulerian momentum equation;
  \item there is no separate \(\rho_0\chi_{tt}\) equation in the coupled model.
\end{itemize}
\end{column}
\end{columns}
\begin{punchbox}
\small Extensions can treat density or viscosity contrasts, but they require additional terms beyond the classical formulation used here.
\end{punchbox}
\end{frame}

\section{Strong form of the coupled problem}

% ----- 5. Coupled equations -----
\begin{frame}{Continuous coupled equations}
\small
\[
\boxed{
  \rho\!\left(
  \frac{\partial\uvec}{\partial t}
  +\,
  (\uvec\!\cdot\!\gradx)\uvec
  \right)
  =-\gradx p+\mu\gradx^2\uvec+\fvec
},
\qquad
\boxed{\divx\uvec=0}.
\]
\[
  \fvec(\xvec,t)
  =
  \int_{B_0}\!\big(\Div\Pten+\Gvec\big)
  \delta\!\left(\xvec-\chi(\Xvec,t)\right)\,d\Xvec
  -
  \int_{\partial B_0}\!\Pten\Nvec\,
  \delta\!\left(\xvec-\chi(\Xvec,t)\right)\,dA .
\]
\[
\boxed{
  \frac{\partial\chi}{\partial t}(\Xvec,t)
  =
  \int_{\Omega}\uvec(\xvec,t)
  \delta\!\left(\xvec-\chi(\Xvec,t)\right)\,d\xvec
  =\uvec(\chi(\Xvec,t),t)
}.
\]
\begin{punchbox}
\small Here \(\Gvec\) is a prescribed force per reference volume. The continuum equations use the singular Dirac distribution \(\delta\); discretization replaces it by \(\delta_h\).
\end{punchbox}
\soft{\scriptsize Griffith and Luo (2017), Eqs.\ (4)--(7), extended here by the body-force term \(\Gvec\).}
\end{frame}

% ----- 6. Force anatomy -----
\begin{frame}{Volume force and transmission force}
\[
  \fvec
  =
  \underbrace{\int_{B_0}(\Div\Pten+\Gvec)\,
  \delta(\xvec-\chi)\,d\Xvec}_{\text{volume contribution}}
  +
  \underbrace{\int_{\partial B_0}(-\Pten\Nvec)\,
  \delta(\xvec-\chi)\,dA}_{\text{transmission contribution}} .
\]
\begin{center}
\renewcommand{\arraystretch}{1.25}
\begin{tabular}{lll}
\toprule
Contribution & Reference support & Character in the fluid \\
\midrule
\(\Div\Pten+\Gvec\) & throughout \(B_0\) & regular body-force density \\
\(-\Pten\Nvec\) & on \(\partial B_0\) & singular interfacial force layer \\
\bottomrule
\end{tabular}
\end{center}
\begin{punchbox}
\small A nonzero transmission force can generate jumps in pressure or viscous stress. Omitting it from the strong form changes the model unless its effect is retained weakly.
\end{punchbox}
\end{frame}

% ----- 7. PK1 -----
\begin{frame}{First Piola--Kirchhoff traction}
\begin{columns}[T,onlytextwidth]
\begin{column}{0.52\textwidth}
For a reference surface with unit normal \(\Nvec\),
\[
  \bm{t}_0=\Pten\Nvec,
  \qquad
  (t_0)_i=P_{iI}N_I .
\]
\begin{itemize}\itemsep5pt
  \item \(i\) is a spatial index.
  \item \(I\) is a material index.
  \item \(\Pten\) is generally nonsymmetric.
\end{itemize}
\end{column}
\begin{column}{0.48\textwidth}
\[
  \Pten=J\sigten^s\Ften^{-\mathsf T},
  \qquad
  \sigten^s=\frac{1}{J}\Pten\Ften^{\mathsf T}.
\]
\[
  [\Pten]=[\bm{t}_0]
  =\mathrm{force/reference\ area}.
\]
\begin{punchbox}
\small \(\Pten\Nvec\) is the traction on the solid. The transmission force \(-\Pten\Nvec\) is not a third-law reaction: it is the singular surface layer that appears because \(\sigten^{e}\) is truncated at \(\partial B_t\) (distributionally, \(\gradx\!\cdot(\sigten^{e}\chi_{B_t})=(\gradx\!\cdot\sigten^{e})\chi_{B_t}-(\sigten^{e}\nvec)\,\delta_{\partial B_t}\), with \(\chi_{B_t}\) the indicator function of \(B_t\)).
\end{punchbox}
\end{column}
\end{columns}
\end{frame}

% ----- 10. Material force balance -----
\begin{frame}{Force balance on a material subregion}
Let \(V_0\subset B_0\) be fixed in reference coordinates. Its contact and body forces are
\[
  \int_{\partial V_0}\Pten\Nvec\,dA,
  \qquad
  \int_{V_0}\Gvec\,d\Xvec .
\]
The reference divergence theorem gives
\[
  \int_{\partial V_0}\Pten\Nvec\,dA
  =
  \int_{V_0}\Div\Pten\,d\Xvec .
\]
Hence the net force on \(V_0\) is
\[
  \int_{V_0}(\Div\Pten+\Gvec)\,d\Xvec .
\]
\begin{punchbox}
\small \(\Div\Pten\) is a spatial force vector per reference volume even though differentiation occurs with respect to material coordinates.
\end{punchbox}
\end{frame}

% ----- 11. Reference momentum bridge -----
\begin{frame}{Reference momentum balance and the IBFE bridge}
For a stand-alone solid, localization of material momentum balance gives
\[
  \rho_0(\Xvec)\,\chi_{tt}
  =\Div\Pten+\Gvec
  \qquad\text{in }B_0.
\]
\begin{columns}[T,onlytextwidth]
\begin{column}{0.49\textwidth}
\textbf{Stand-alone solid}

The Lagrangian equation carries solid inertia through \(\rho_0\chi_{tt}\).
\end{column}
\begin{column}{0.51\textwidth}
\textbf{Neutrally buoyant IBFE}

The common Eulerian momentum equation carries fluid and solid inertia. The structural model supplies force, not a second inertia equation.
\end{column}
\end{columns}
\begin{punchbox}
\small IBFE spreads \(\Div\Pten+\Gvec\) and the boundary reaction \(-\Pten\Nvec\) to the Eulerian grid.
\end{punchbox}
\end{frame}

\section{Weak form of the structural force}

% ----- 12. Why weak -----
\begin{frame}{Why IBFE uses a weak structural force}
\begin{itemize}\itemsep8pt
  \item \(C^0\) finite elements give a piecewise stress, so \(\Div\Pten\) need not exist pointwise across element boundaries.
  \item Virtual work needs only first derivatives of the motion and the test function.
  \item The boundary transmission force remains part of the formulation.
  \item The result is a Lagrangian force field ready for Eulerian spreading.
\end{itemize}
\begin{punchbox}
\small The weak form is the structural-force definition used by the numerical method, not merely an alternative notation for a smooth strong solution.
\end{punchbox}
\end{frame}

% ----- 13. Test and contraction -----
\begin{frame}{Test functions and double contraction}
Choose an arbitrary smooth Lagrangian test function \(\Vvec(\Xvec)\). For \(\Pten,\Grad\Vvec\in\mathbb{R}^{d\times d}\),
\[
  \Pten:\Grad\Vvec
  =\operatorname{tr}\!\left(\Pten^{\mathsf T}\Grad\Vvec\right)
  =\sum_{i=1}^{d}\sum_{I=1}^{d}
  P_{iI}\frac{\partial V_i}{\partial X_I}.
\]
Multiplying the interior force by \(\Vvec\) and integrating gives
\[
  \int_{B_0}(\Div\Pten+\Gvec)\cdot\Vvec\,d\Xvec .
\]
\begin{punchbox}
\small The contraction pairs a spatial force index \(i\) with a material derivative index \(I\), consistent with the two-point character of \(\Pten\).
\end{punchbox}
\end{frame}

% ----- 14. Integration by parts -----
\begin{frame}{Integration by parts retains the boundary term}
The tensor product rule is
\[
  \Div(\Pten^{\mathsf T}\Vvec)
  =
  (\Div\Pten)\cdot\Vvec
  +\Pten:\Grad\Vvec .
\]
After integration over \(B_0\),
\[
  \int_{B_0}(\Div\Pten)\cdot\Vvec\,d\Xvec
  =
  -\int_{B_0}\Pten:\Grad\Vvec\,d\Xvec
  +\int_{\partial B_0}(\Pten\Nvec)\cdot\Vvec\,dA .
\]
\begin{punchbox}
\small No traction-free assumption is made. The next two formulations account for the reaction \(-\Pten\Nvec\) either implicitly in one weak density or explicitly as a surface density.
\end{punchbox}
\end{frame}

% ----- 15. Unified weak form -----
\begin{frame}{Unified weak form}
Define one Lagrangian force density \(\Fvec\) by
\[
\boxed{
  \int_{B_0}\Fvec\cdot\Vvec\,d\Xvec
  =
  -\int_{B_0}\Pten:\Grad\Vvec\,d\Xvec
  +\int_{B_0}\Gvec\cdot\Vvec\,d\Xvec
  \quad\text{for every }\Vvec
}.
\]
The test functions are \emph{not} required to vanish on \(\partial B_0\). Integration by parts gives
\[
  \int_{B_0}\Fvec\cdot\Vvec\,d\Xvec
  =
  \int_{B_0}(\Div\Pten+\Gvec)\cdot\Vvec\,d\Xvec
  -\int_{\partial B_0}(\Pten\Nvec)\cdot\Vvec\,dA .
\]
\begin{punchbox}
\small The single weak density \(\Fvec\) packages the interior force and the boundary transmission force. Independently prescribed external tractions would still contribute separate boundary work.
\end{punchbox}
\soft{\scriptsize Griffith and Luo (2017), Eqs.\ (10)--(14).}
\end{frame}

% ----- 9. Partitioned weak form -----
\begin{frame}{Partitioned weak form}
Keep the volume and boundary forces separate. Find \(\Fvec_{\mathrm{vol}}\) such that
\[
  \int_{B_0}\Fvec_{\mathrm{vol}}\cdot\Vvec\,d\Xvec
  =
  -\int_{B_0}\Pten:\Grad\Vvec\,d\Xvec
  +\int_{\partial B_0}(\Pten\Nvec)\cdot\Vvec\,dA
  +\int_{B_0}\Gvec\cdot\Vvec\,d\Xvec .
\]
\par\medskip\noindent Then define
\[
  \Tvec=-\Pten\Nvec
  \quad\text{on }\partial B_0,
\]
and spread the two fields independently:
\[
  \fvec
  =
  \int_{B_0}\Fvec_{\mathrm{vol}}\,\delta(\xvec-\chi)\,d\Xvec
  +
  \int_{\partial B_0}\Tvec\,\delta(\xvec-\chi)\,dA .
\]
\soft{\scriptsize Griffith and Luo (2017), Eqs.\ (15)--(21).}
\end{frame}

% ----- 10. Compare -----
\begin{frame}{Unified and partitioned formulations}
\begin{columns}[T,onlytextwidth]
\begin{column}{0.49\textwidth}
\textbf{Unified}
\begin{itemize}\itemsep6pt
  \item one total Lagrangian density \(\Fvec\);
  \item boundary transmission represented through its weak definition;
  \item compact assembly and one spreading operator.
\end{itemize}
\end{column}
\begin{column}{0.51\textwidth}
\textbf{Partitioned}
\begin{itemize}\itemsep6pt
  \item volume density \(\Fvec_{\mathrm{vol}}\) plus surface density \(\Tvec\);
  \item transmission force spread directly from \(\partial B_0\);
  \item can improve accuracy for nonzero traction and a relatively coarse FE mesh.
\end{itemize}
\end{column}
\end{columns}
\medskip
\begin{punchbox}
\small The formulations are equivalent in the continuum. Their finite-dimensional approximations differ, so neither discrete form is universally superior.
\end{punchbox}
\end{frame}

\section{Discretization and coupling operators}

% ----- 18. Continuous cycle -----
\begin{frame}{Complete continuous coupling cycle}
\begin{enumerate}\itemsep6pt
  \item Evaluate the structural stress \(\Pten(\Ften)\).
  \item Define \(\Fvec\) by virtual work:
  \[
    \int_{B_0}\Fvec\cdot\Vvec\,d\Xvec
    =-\int_{B_0}\Pten:\Grad\Vvec\,d\Xvec
    +\int_{B_0}\Gvec\cdot\Vvec\,d\Xvec .
  \]
  \item Spread force and solve the incompressible fluid equations:
  \[
    \fvec(\xvec,t)=\int_{B_0}\Fvec(\Xvec,t)\,
    \delta(\xvec-\chi(\Xvec,t))\,d\Xvec .
  \]
  \item Move the structure with \(\chi_t=\uvec\circ\chi\).
\end{enumerate}
\begin{punchbox}
\small Spatial discretization replaces \(\delta\) with \(\delta_h\) and evaluates the interaction integrals by quadrature.
\end{punchbox}
\end{frame}

% ----- 19. FE fields -----
\begin{frame}{Finite-element motion and force fields}
\small
With basis functions \(\{\phi_\ell\}_{\ell=1}^{M}\), approximate the motion and unified force by
\[
  \chi_h(\Xvec,t)=\sum_{\ell=1}^{M}\chi_\ell(t)\phi_\ell(\Xvec),
  \qquad
  \Fvec_h(\Xvec,t)=\sum_{\ell=1}^{M}\Fvec_\ell(t)\phi_\ell(\Xvec).
\]
\[
  \Ften^{\,h}=\Grad\chi_h,
  \qquad
  \Pten_h=\Pten(\Ften^{\,h}).
\]
\begin{itemize}\itemsep6pt
  \item \(\chi_h\) is continuous for standard \(C^0\) elements.
  \item \(\Ften^{\,h}\) and \(\Pten_h\) are generally discontinuous across element boundaries.
  \item The weak formulation requires elementwise first derivatives, not a pointwise \(\Div\Pten_h\).
\end{itemize}
\begin{punchbox}
\small The coefficients \(\chi_\ell\) carry geometry. The coefficients \(\Fvec_\ell\) represent the projected unified Lagrangian force.
\end{punchbox}
\end{frame}

% ----- 20. Force projection -----
\begin{frame}{Mass-matrix projection of the weak force}
Testing the unified weak definition with \(\Vvec=\phi_m\bm{e}_i\) yields
\[
  \sum_{\ell=1}^{M}M_{m\ell}\Fvec_\ell=\bm{B}_m,
  \qquad
  M_{m\ell}=\int_{B_0}\phi_m\phi_\ell\,d\Xvec,
\]
\[
  \bm{B}_m
  =
  -\int_{B_0}\Pten_h\Grad\phi_m\,d\Xvec
  +\int_{B_0}\Gvec\phi_m\,d\Xvec .
\]
In block-vector form,
\[
  M\Fcoef=\bm{B},
  \qquad
  \Fcoef=
  \begin{bmatrix}
  \Fvec_1&\cdots&\Fvec_M
  \end{bmatrix}^{\mathsf T}.
\]
\begin{punchbox}
\small This \(L^2\) projection converts virtual work into the nodal coefficient vector used by the coupling operators.
\end{punchbox}
\end{frame}

% ----- 12. Quadrature coupling -----
\begin{frame}{Quadrature-point force spreading}
Let \(\Xvec_Q^e\) and \(w_Q^e\) be quadrature points and weights in element \(e\). The discrete force-spreading operator has the form
\[
  (\fvec_h)_i
  =
  \sum_e\sum_Q
  \Fvec_h(\Xvec_Q^e)\,
  \delta_h\!\left(
  \xvec_i-\chi_h(\Xvec_Q^e)
  \right)w_Q^e
  =(\mathcal{S}_h\Fvec_h)_i .
\]
\begin{itemize}\itemsep5pt
  \item The FE nodes determine \(\chi_h\) and the stress approximation.
  \item Interaction values are evaluated at quadrature points inside each element.
  \item Adaptive quadrature can keep the interaction net dense as the solid deforms.
\end{itemize}
\begin{punchbox}
\small The FE mesh can therefore be coarser than the Eulerian grid without leaving gaps in the coupling stencil, provided the quadrature-point net is sufficiently dense.
\end{punchbox}
\soft{\scriptsize Griffith and Luo (2017), Eqs.\ (31)--(33) and Fig.\ 2.}
\end{frame}

% ----- 13. Adjoint restriction -----
\begin{frame}{Adjoint velocity restriction}
Let \(M\) be the Lagrangian mass matrix and \(H\) the Eulerian grid inner-product matrix. Define the nodal velocity by
\[
  \boxed{
  \dot{\bm{\chi}}
  =
  \mathcal{R}_h\uvec_h
  =
  M^{-1}\mathcal{S}_h^{\mathsf T}H\uvec_h
  }.
\]
Let \(\Fcoef\) denote the block vector of Lagrangian force coefficients. Then the discrete power identity holds:
\[
  \underbrace{\Fcoef^{\mathsf T}M\dot{\bm{\chi}}}_{\text{Lagrangian force power}}
  =
  \underbrace{(\mathcal{S}_h\Fcoef)^{\mathsf T}H\uvec_h}_{\text{Eulerian force power}} .
\]
\begin{punchbox}
\small Force spreading and velocity restriction are adjoint under their discrete inner products. This prevents the coupling step from creating or destroying energy by itself.
\end{punchbox}
\soft{\scriptsize Griffith and Luo (2017), Eqs.\ (35)--(40); Peskin (2002).}
\end{frame}

% ----- 14. Incompressibility -----
\begin{frame}{Incompressibility of the immersed material}
Because \(\chi_t=\uvec\circ\chi\), Jacobi's formula gives
\[
  \frac{\partial J}{\partial t}
  =
  J\,\operatorname{tr}
  \!\left(\Ften^{-1}\frac{\partial\Ften}{\partial t}\right)
  =
  J\,(\divx\uvec)\circ\chi
  =0.
\]
Therefore
\[
  J(\Xvec,t)=J(\Xvec,0),
  \qquad
  J\equiv1\ \text{only when }J(\Xvec,0)=1.
\]
\begin{columns}[T,onlytextwidth]
\begin{column}{0.53\textwidth}
At the continuum level, no slip and an incompressible Eulerian velocity preserve solid volume.
\end{column}
\begin{column}{0.47\textwidth}
At the discrete level, monitor \(J_h\) and use an appropriate volumetric stabilization when needed.
\end{column}
\end{columns}
\begin{punchbox}
\small In \(d\) dimensions, a common stabilized construction uses \(\overline{\Ften}=J^{-1/d}\Ften\) for the isochoric response plus a separate penalty \(U(J)\).
\end{punchbox}
\soft{\scriptsize Griffith and Patankar (2020), Sec.\ 2.2.2.}
\end{frame}

\section{Active stress and biological examples}

% ----- 15. Active split -----
\begin{frame}{Passive and active stress}
Biological solids may generate active stress in addition to passive elastic resistance:
\[
  \boxed{\Pten=\Pten^{p}+\Pten^{a}},
  \qquad
  \Pten^{p}=\frac{\partial W}{\partial\Ften}.
\]
The active term \(\Pten^a\) may depend on fiber direction, an activation variable, and time. It need not derive from a stored-energy density.
\begin{columns}[T,onlytextwidth]
\begin{column}{0.49\textwidth}
\textbf{Passive response}

Stores and releases elastic energy through \(W(\Ften)\).
\end{column}
\begin{column}{0.51\textwidth}
\textbf{Active response}

Supplies mechanical power through prescribed or regulated contraction.
\end{column}
\end{columns}
\begin{punchbox}
\small Hereafter \(\Pten\) denotes the total structural PK1 stress. Setting \(\Pten^a=\bm{0}\) recovers the passive convention used in the solid-mechanics module.
\end{punchbox}
\end{frame}

% ----- 25. Active assembly -----
\begin{frame}{Active stress in the weak-force assembly}
Substituting \(\Pten=\Pten^p+\Pten^a\) into the unified weak form gives
\[
  \int_{B_0}\Fvec\cdot\Vvec\,d\Xvec
  =
  -\int_{B_0}
  \left(\Pten^p+\Pten^a\right):\Grad\Vvec\,d\Xvec
  +\int_{B_0}\Gvec\cdot\Vvec\,d\Xvec .
\]
At each structural quadrature point:
\begin{enumerate}\itemsep5pt
  \item evaluate \(\Pten^p\) from the constitutive law;
  \item evaluate \(\Pten^a\) from the activation model;
  \item add the stresses before assembling the force projection.
\end{enumerate}
\begin{punchbox}
\small Once expressed as PK1 stress, passive and active contributions use the same FE assembly and fluid-coupling operators.
\end{punchbox}
\end{frame}

% ----- 16. Fiber stress -----
\begin{frame}{Fiber-oriented active stress}
Let \(\fzero(\Xvec)\) be a unit reference fiber direction and define
\[
  \avec=\Ften\fzero,
  \qquad
  \lambda_f=\|\avec\|.
\]
One unnormalized mapped-fiber convention is
\[
  \sigten^a=T\,\avec\otimes\avec,
\]
which pulls back to
\[
\boxed{
  \Pten^a
  =
  J\sigten^a\Ften^{-\mathsf T}
  =
  JT\,\avec\otimes\fzero
  =
  JT\,\Ften\fzero\fzero^{\mathsf T}
}.
\]
\begin{punchbox}
\small This convention uses the mapped fiber \(\avec\), not its unit direction. A model based on \(\mvec=\avec/\lambda_f\) contains different powers of \(\lambda_f\); the scalar \(T\) must be calibrated for the convention used.
\end{punchbox}
\end{frame}

% ----- 27. Jellyfish fibers -----
\begin{frame}{Jellyfish circumferential reference fibers}
The subumbrellar swimming musculature is represented by the reference direction
\[
  \fzero(\Xvec)
  =
  \left(-\frac{X_2}{R},\,\frac{X_1}{R},\,0\right),
  \qquad
  R=\sqrt{X_1^2+X_2^2}>0,
\]
where \(X_1,X_2,X_3\) are the scalar components of the reference point \(\Xvec\) and the bell axis is \(X_3\).
With the mapped-fiber convention, the active PK1 stress is
\[
  \Pten^a
  =
  JT(\Xvec,t)\,\Ften\fzero\fzero^{\mathsf T}.
\]
\begin{itemize}\itemsep4pt
\small
  \item The direction is tangent to circles centered on the bell axis.
  \item Positive activation produces circumferential contractile stress.
  \item At \(R=0\), the circumferential direction is undefined. Avoid axis quadrature points, regularize the field, or set the activation to zero there.
\end{itemize}
\soft{\scriptsize Active-stress model: convention and activation form as used in Hoover, Griffith \& Miller (2017); check against that paper before reusing the constants.}
\end{frame}

% ----- 28. Jellyfish activation -----
\begin{frame}{Jellyfish activation field}
Write the tension magnitude as
\[
  T(\Xvec,t)
  =
  T_{\max}\,\gamma(\Xvec,t),
  \qquad
  \gamma(\Xvec,t)=\alpha(X_3)\,\beta(t),
\]
\[
  0\le\alpha,\beta,\gamma\le1.
\]
\begin{columns}[T,onlytextwidth]
\begin{column}{0.49\textwidth}
\textbf{Spatial profile \(\alpha(X_3)\)}

Selects the active band of the bell. A smooth transition avoids a discontinuous stress field.
\end{column}
\begin{column}{0.51\textwidth}
\textbf{Temporal profile \(\beta(t)\)}

Controls contraction and relaxation over the pulse cycle. Smooth ramps reduce impulsive forcing.
\end{column}
\end{columns}
\begin{punchbox}
\small Report \(T_{\max}\), the spatial and temporal profiles, and the fiber-stress convention together. A value calibrated for a unit-current-fiber convention cannot be transferred unchanged.
\end{punchbox}
\soft{\scriptsize Active-stress example: Hoover, Griffith, and Miller (2017). Convention and activation form as used in that paper; check against it before reusing the constants.}
\end{frame}

% ----- 18. Tether -----
\begin{frame}{Localized tether as an external body force}
An optional restoring force can anchor a selected region:
\[
  \Gvec(\Xvec,t)=
  \begin{cases}
  \kappa\,[\chi(\Xvec,0)-\chi(\Xvec,t)],
    &R\le R_0,\\[3pt]
  \bm{0},&R>R_0 .
  \end{cases}
\]
\begin{itemize}\itemsep6pt
  \item \(\Gvec\) is a force per reference volume.
  \item It enters the unified weak form through
  \(\int_{B_0}\Gvec\cdot\Vvec\,d\Xvec\).
  \item \(\kappa\) has units of force per reference volume per displacement.
\end{itemize}
\begin{punchbox}
\small A large tether stiffness can dominate the explicit structural time-step restriction. Verify that the tether does not alter the intended deformation outside the anchored region.
\end{punchbox}
\end{frame}

\section{Worked example: uniaxial stretch (prescribed axial deformation)}

% ----- 30. Stretch problem -----
\begin{frame}{Homogeneous-stretch problem and test types}
For a block \(B_0=[0,L]\times[0,H]\times[0,W]\), prescribe
\[
  \chi(\Xvec)
  =
  \left(
  \lambda X_1,\,
  \lambda_{\perp}X_2,\,
  \lambda_{\perp}X_3
  \right).
\]
\begin{itemize}\itemsep6pt
  \item \(\lambda_\perp=1\): uniaxial stretch (prescribed axial deformation);
  \item \(\lambda_\perp=\lambda^{-1/2}\): prescribed isochoric stretch;
  \item traction-free lateral faces: solve for \(\lambda_\perp\) from the constitutive law.
\end{itemize}
\begin{punchbox}
\small Prescribing the lateral stretch and prescribing zero lateral traction are different experiments.
\end{punchbox}
\end{frame}

% ----- 31. Stretch kinematics -----
\begin{frame}{Homogeneous-stretch kinematics}
For the prescribed motion,
\[
  \Ften=\Grad\chi
  =\operatorname{diag}(\lambda,\lambda_\perp,\lambda_\perp),
  \qquad
  J=\lambda\lambda_\perp^2.
\]
Line and volume elements transform according to
\[
  d\xvec=\Ften\,d\Xvec,
  \qquad
  dv=J\,dV.
\]
\begin{punchbox}
\small The motion preserves orientation only for \(\lambda>0\) and \(\lambda_\perp>0\). The special choice \(\lambda_\perp=\lambda^{-1/2}\) gives \(J=1\).
\end{punchbox}
\end{frame}

% ----- 32. Stretch strain -----
\begin{frame}{Strain measures and invariants for the stretch}
\[
  \Cten=\Ften^{\mathsf T}\Ften
  =\operatorname{diag}
  (\lambda^2,\lambda_\perp^2,\lambda_\perp^2),
\]
\[
  \Eten=\tfrac12(\Cten-\Identity)
  =\tfrac12\operatorname{diag}
  (\lambda^2-1,\lambda_\perp^2-1,\lambda_\perp^2-1).
\]
The invariants needed below are
\[
  I_1=\lambda^2+2\lambda_\perp^2,
  \qquad
  I_3=\det\Cten=J^2=\lambda^2\lambda_\perp^4.
\]
\end{frame}

% ----- 33. Neo-Hookean energy -----
\begin{frame}{Compressible neo-Hookean energy and parameters}
Consider the energy per reference volume
\[
  W(\Ften)
  =
  \frac{\mu_s}{2}(I_1-3)
  -\mu_s\ln J
  +\frac{\lambda_{\mathrm L}}{2}(\ln J)^2,
  \qquad
  I_1=\operatorname{tr}(\Ften^{\mathsf T}\Ften).
\]
\begin{columns}[T,onlytextwidth]
\begin{column}{0.52\textwidth}
\(\mu_s\) is the infinitesimal shear modulus.
\end{column}
\begin{column}{0.48\textwidth}
\(\lambda_{\mathrm L}\) is the first Lamé parameter, not the bulk modulus.
\end{column}
\end{columns}
\begin{punchbox}
\small In the small-strain limit, the bulk modulus is \(K=\lambda_{\mathrm L}+2\mu_s/3\).
\end{punchbox}
\end{frame}

% ----- 34. Differentiate energy -----
\begin{frame}{Differentiating the neo-Hookean energy}
Use
\[
  \frac{\partial I_1}{\partial\Ften}=2\Ften,
  \qquad
  \frac{\partial\ln J}{\partial\Ften}=\Ften^{-\mathsf T}.
\]
Then
\[
\boxed{
  \Pten
  =
  \mu_s\left(\Ften-\Ften^{-\mathsf T}\right)
  +\lambda_{\mathrm L}(\ln J)\Ften^{-\mathsf T}
}.
\]
\begin{punchbox}
\small At \(\Ften=\Identity\), \(J=1\) and \(\Pten=\bm{0}\), so the reference configuration is stress free.
\end{punchbox}
\end{frame}

% ----- 20. Stretch stress -----
\begin{frame}{Stress response of the homogeneous stretch}
For \(\Ften=\operatorname{diag}(\lambda,\lambda_\perp,\lambda_\perp)\),
\[
  P_{11}
  =
  \mu_s\!\left(\lambda-\frac{1}{\lambda}\right)
  +\lambda_{\mathrm L}\frac{\ln J}{\lambda},
\]
\[
  P_{22}=P_{33}
  =
  \mu_s\!\left(\lambda_\perp-\frac{1}{\lambda_\perp}\right)
  +\lambda_{\mathrm L}\frac{\ln J}{\lambda_\perp},
  \qquad
  J=\lambda\lambda_\perp^2 .
\]
\begin{itemize}\itemsep5pt
  \item If \(\lambda_\perp=1\), then \(P_{22}=P_{33}=\lambda_{\mathrm L}\ln\lambda\), generally nonzero.
  \item A uniaxial-stress test requires \(P_{22}=P_{33}=0\) and determines \(\lambda_\perp\).
  \item An incompressibility constraint \(J=1\) introduces an undetermined pressure and is a different constitutive problem.
\end{itemize}
\end{frame}

% ----- 21. Boundary forces in test -----
\begin{frame}{Uniform stress acts through the boundary}
For the homogeneous test, \(\Pten\) is spatially constant, so
\[
  \Div\Pten=\bm{0}
  \quad\text{inside }B_0.
\]
The unified virtual-work term nevertheless remains nonzero:
\[
  -\int_{B_0}\Pten:\Grad\Vvec\,d\Xvec
  =
  -\int_{\partial B_0}(\Pten\Nvec)\cdot\Vvec\,dA .
\]
In the partitioned description,
\[
  \Fvec_{\mathrm{vol}}=\bm{0},
  \qquad
  \Tvec=-\Pten\Nvec .
\]
\begin{punchbox}
\small Constant stress does not mean zero structural force. It means zero interior divergence, with the load carried by boundary tractions.
\end{punchbox}
\end{frame}

% ----- 38. Kinematic versus coupled -----
\begin{frame}{Kinematic test and fully coupled IBFE simulation}
\begin{columns}[T,onlytextwidth]
\begin{column}{0.49\textwidth}
\textbf{Kinematic constitutive test}
\begin{itemize}\itemsep6pt
  \item prescribe \(\chi\);
  \item compute \(\Ften,J,\Pten\);
  \item compare against closed-form values;
  \item use the result to verify material routines.
\end{itemize}
\end{column}
\begin{column}{0.51\textwidth}
\textbf{Fully coupled simulation}
\begin{itemize}\itemsep6pt
  \item assemble and spread \(\Fvec\);
  \item solve for \(\uvec\) and \(p\);
  \item restrict \(\uvec\) to update \(\chi\);
  \item repeat in time.
\end{itemize}
\end{column}
\end{columns}
\begin{punchbox}
\small The homogeneous stretch validates the solid-side formulas. It is not itself a dynamically coupled IBFE solution.
\end{punchbox}
\end{frame}

\section{Verification and references}

% ----- 22. Verification -----
\begin{frame}{IBFE verification checklist}
\begin{columns}[T,onlytextwidth]
\begin{column}{0.51\textwidth}
\textbf{Solid mechanics}
\begin{itemize}\itemsep4pt
  \item verify \(J>0\) and monitor volume drift;
  \item test objectivity with a rigid rotation;
  \item compare \(\Pten\) with a closed-form homogeneous deformation;
  \item check boundary tractions and sign conventions.
\end{itemize}
\end{column}
\begin{column}{0.49\textwidth}
\textbf{Eulerian coupling}
\begin{itemize}\itemsep4pt
  \item verify total-force consistency;
  \item test the discrete power identity;
  \item refine the grid, FE mesh, quadrature net, and time step separately;
  \item check sensitivity to the kernel and volumetric stabilization.
\end{itemize}
\end{column}
\end{columns}
\begin{punchbox}
\centering
\small Constitutive accuracy, FE assembly, and Lagrangian--Eulerian coupling require separate tests.
\end{punchbox}
\end{frame}

% ----- 23. References -----
\begin{frame}{References and resources}
\footnotesize
\begin{itemize}\itemsep4pt
  \item B. E. Griffith and X. Luo, \emph{Hybrid finite difference/finite element immersed boundary method}, \emph{Int. J. Numer. Meth. Biomed. Engng.} 33 (2017), e2888, \href{https://doi.org/10.1002/cnm.2888}{doi:10.1002/cnm.2888}.
  \item D. Boffi, L. Gastaldi, L. Heltai, and C. S. Peskin, \emph{On the hyper-elastic formulation of the immersed boundary method}, \emph{Comput. Methods Appl. Mech. Engrg.} 197 (2008), 2210--2231.
  \item B. E. Griffith and N. A. Patankar, \emph{Immersed methods for fluid--structure interaction}, \emph{Annu. Rev. Fluid Mech.} 52 (2020), 421--448, \href{https://doi.org/10.1146/annurev-fluid-010719-060228}{doi:10.1146/annurev-fluid-010719-060228}.
  \item C. S. Peskin, \emph{The immersed boundary method}, \emph{Acta Numerica} 11 (2002), 479--517, \href{https://doi.org/10.1017/S0962492902000077}{doi:10.1017/S0962492902000077}.
  \item A. P. Hoover, B. E. Griffith, and L. A. Miller, \emph{Quantifying performance in the medusan mechanospace with an actively swimming three-dimensional jellyfish model}, \emph{J. Fluid Mech.} 813 (2017), 1112--1155, \href{https://doi.org/10.1017/jfm.2017.3}{doi:10.1017/jfm.2017.3}.
\end{itemize}
\begin{punchbox}
\centering
\small IBAMR documentation and source: \href{https://ibamr.github.io}{ibamr.github.io}
\end{punchbox}
\end{frame}

\end{document}
